%
% ggz.tex -- eukldische Eigenschaft für ganze gaussche Zahlen
%
% (c) 2021 Prof Dr Andreas Müller, OST Ostschweizer Fachhochschule
%
\documentclass[tikz]{standalone}
\usepackage{amsmath}
\usepackage{times}
\usepackage{txfonts}
\usepackage{pgfplots}
\usepackage{csvsimple}
\usetikzlibrary{arrows,intersections,math}
\definecolor{darkred}{rgb}{0.8,0,0}
\definecolor{darkgreen}{rgb}{0,0.6,0}
\begin{document}
\def\dx{0.5}
\def\dy{0.5}
\def\skala{1}
\def\bx{4}
\def\by{3}
\pgfmathparse{sqrt(\bx*\bx+\by*\by)}
\xdef\r{\pgfmathresult}
\def\ax{11}
\def\ay{7}
\def\gausspunkt#1#2{ ({(#1)*\dx},{(#2)*\dy}) }
\def\gitterpunkt#1#2{ ({((#1)*\bx-(#2)*\by)*\dx},{((#1)*\by+(#2)*\bx)*\dy}) }
\def\kreis#1#2{
	\fill[color=darkgreen!20,opacity=0.5]
		\gitterpunkt{#1}{#2} circle[radius={\r*\dx}];
}
\def\kreisrand#1#2{
	\draw[color=darkgreen]
		\gitterpunkt{#1}{#2} circle[radius={\r*\dx}];
}
\begin{tikzpicture}[>=latex,thick,scale=\skala]
\clip (-3.1,-3.1) rectangle +(12.6,9.6);

\draw[->] (-3.1,0) -- (9.3,0) coordinate[label={$\operatorname{Re}$}];
\draw[->] (0,-3.1) -- (0,6.3) coordinate[label={right:$\operatorname{Im}$}];

\foreach \t in {-6,...,18}{
	\draw[line width=0.2pt] ({\t*\dx},-3.1) -- ++(0,9.2);
}
\foreach \t in {-6,...,12}{
	\draw[line width=0.2pt] (-3.1,{\t*\dy}) -- ++(12.2,0);
}
\foreach \x in {-6,...,18}{
	\foreach \y in {-6,...,12}{
		\fill ({\x*\dx},{\y*\dy}) circle[radius=0.03];
	}
}

\fill[color=darkred!20,opacity=0.2]
	\gitterpunkt{2}{-1} --
	\gitterpunkt{3}{-1} --
	\gitterpunkt{3}{0} --
	\gitterpunkt{2}{0} -- cycle;

\begin{scope}
	\clip
	\gitterpunkt{2}{-1} --
	\gitterpunkt{3}{-1} --
	\gitterpunkt{3}{0} --
	\gitterpunkt{2}{0} -- cycle;
	\kreis{2}{-1}
	\kreis{3}{-1}
	\kreis{3}{0}
	\kreis{2}{0}
	\kreisrand{2}{-1}
	\kreisrand{3}{-1}
	\kreisrand{3}{0}
	\kreisrand{2}{0}
\end{scope}

\draw[->,color=darkred,line width=1.2pt] (0,0) -- \gitterpunkt{1}{0};
\draw[->,color=darkred,line width=1.2pt] (0,0) -- \gitterpunkt{0}{1};
\node[color=darkred] at \gitterpunkt{0.5}{0} [below right] {$b$};
\node[color=darkred] at \gitterpunkt{0}{0.5} [below left] {$ib$};

\draw[->,color=darkred,line width=1.2pt]
	\gitterpunkt{2}{-1} -- \gitterpunkt{3}{-1};
\draw[->,color=darkred,line width=1.2pt]
	\gitterpunkt{2}{-1} -- \gitterpunkt{2}{0};
\node[color=darkred] at \gitterpunkt{2.5}{-1} [below right] {$b$};
\node[color=darkred] at \gitterpunkt{2}{-0.5} [below left] {$ib$};

\draw[color=orange] \gitterpunkt{3}{0} -- \gausspunkt{\ax}{\ay};
\node[color=orange] at \gausspunkt{\ax+0.7}{\ay+0.9} {$r$};

\begin{scope}
	\clip (-3.1,-3.1) rectangle (9.1,6.1);
	\foreach \x in {-4,...,4}{
		\foreach \y in {-5,...,5}{
		\fill[color=darkred] \gitterpunkt{\x}{\y} circle[radius=0.05];
		}
	}
	\foreach \x in {-4,...,4}{
		\draw[color=darkred] \gitterpunkt{\x}{-5}
			-- ++\gitterpunkt{0}{10};
	}
	\foreach \y in {-4,...,4}{
		\draw[color=darkred] \gitterpunkt{-5}{\y}
			-- ++ \gitterpunkt{10}{0};
	}
\end{scope}

\fill[color=blue] \gausspunkt{\ax}{\ay} circle[radius=0.08];
\node[color=blue] at \gausspunkt{\ax}{\ay} [below] {$a$};

\node[color=darkred] at \gitterpunkt{3}{0} [above] {$qb$};
\fill[color=darkred] \gitterpunkt{3}{0} circle[radius=0.08];

\end{tikzpicture}
\end{document}

